\chapter{Shoulder Replacement}

\section*{Introduction \& Clinical Indications}
Shoulder replacement, formally known as glenohumeral arthroplasty, is a sophisticated surgical procedure meticulously designed to alleviate severe, debilitating pain and restore functional mobility in a compromised shoulder joint. This intervention involves the precise excision of diseased or injured articular surfaces of the humeral head and/or the glenoid cavity, which are then meticulously replaced with artificial prosthetic components. The procedure primarily targets conditions affecting the glenohumeral joint, the primary ball-and-socket articulation of the shoulder, which is crucial for a wide range of arm movements. Depending on the patient's specific pathology, particularly the integrity and function of the rotator cuff, the surgeon may elect to perform an anatomic total shoulder arthroplasty (TSA) or a reverse total shoulder arthroplasty (rTSA), each employing distinct biomechanical principles and having specific clinical indications to optimize patient outcomes.

\begin{itemize}
  \item Total Shoulder Arthroplasty (TSA)
  \item Reverse Total Shoulder Arthroplasty (rTSA)
  \item Hemiarthroplasty of the Shoulder
  \item Glenohumeral Arthroplasty
  \item Shoulder Joint Replacement
  \item Arthroplasty of the Shoulder
\end{itemize}

The fundamental principle of shoulder replacement is to meticulously remove damaged articular cartilage and subchondral bone, replacing these surfaces with smooth, biocompatible prosthetic implants to restore pain-free motion and stability. The rationale is to eliminate the source of pain, which is typically bone-on-bone articulation or severe inflammatory changes, and to reconstruct a functional joint. In an anatomic total shoulder arthroplasty (TSA), the procedure aims to precisely replicate the natural anatomy and biomechanics of the glenohumeral joint. This involves replacing the humeral head with a metal ball component, typically affixed to a stem inserted into the humerus, and resurfacing the glenoid (shoulder socket) with a polyethylene component cemented into place. This approach critically relies on an intact and functional rotator cuff to power and stabilize the joint, as the rotator cuff muscles are responsible for fine motor control and elevation.

Conversely, a reverse total shoulder arthroplasty (rTSA) fundamentally alters the biomechanics of the shoulder to compensate for a deficient rotator cuff. In this innovative design, the "ball" component (glenosphere) is fixed to the glenoid, and the "socket" component (humeral cup) is attached to the humerus. This reversal shifts the center of rotation medially and inferiorly, significantly increasing the deltoid muscle's lever arm. This biomechanical alteration allows the deltoid muscle to effectively elevate the arm, even in the complete absence of a functional rotator cuff. The primary technical goals for both procedures are consistent: to eliminate chronic pain, improve the patient's active and passive range of motion, and ultimately enhance the patient's ability to perform essential activities of daily living, thereby significantly improving their quality of life.

The concept of shoulder replacement surgery has a rich history, with early attempts dating back to the late 19th century. The first documented attempt at shoulder arthroplasty was performed by Jules Emile Pean in France in 1893, who implanted a rudimentary platinum and rubber prosthesis in a patient suffering from tuberculosis of the shoulder. While this pioneering effort was not widely adopted and had limited success, it laid the conceptual groundwork for future developments. Modern shoulder arthroplasty truly began to evolve in the 1950s with the groundbreaking work of Dr. Charles Neer II in the United States. In 1955, Neer developed the first widely successful shoulder prosthesis, initially designed for the treatment of complex proximal humerus fractures, and subsequently refined it for arthritic conditions. His innovative designs, particularly the modular humeral component, became the gold standard for decades, establishing the principles of anatomic shoulder replacement.

The development of the reverse total shoulder arthroplasty is a more recent, yet equally transformative, advancement, primarily attributed to Dr. Paul Grammont in France in the 1980s. Grammont's revolutionary design, which reversed the traditional ball-and-socket configuration, fundamentally altered the biomechanics of the shoulder. This innovation revolutionized the treatment of rotator cuff tear arthropathy by enabling the deltoid muscle to effectively compensate for a massive, irreparable rotator cuff tear, thereby restoring active elevation. Subsequent advancements in both anatomic and reverse arthroplasty have focused on improved implant materials (e.g., highly cross-linked polyethylene), enhanced fixation techniques (cemented versus uncemented, porous coatings), and refined surgical approaches, leading to the highly successful and specialized procedures performed today, with ongoing research into patient-specific instrumentation and augmented glenoids.

Shoulder replacement is typically indicated for patients experiencing severe, debilitating shoulder pain and significant functional limitations that have proven refractory to comprehensive conservative treatments. Specific clinical indications include:
\begin{itemize}
  \item Severe primary osteoarthritis of the glenohumeral joint, characterized by extensive articular cartilage loss, subchondral bone sclerosis, and osteophyte formation, leading to bone-on-bone articulation.
  \item Rheumatoid arthritis and other systemic inflammatory arthropathies (e.g., psoriatic arthritis) causing progressive and irreversible joint destruction, pain, and stiffness.
  \item Post-traumatic arthritis resulting from previous shoulder injuries, such as fractures of the humeral head or glenoid, or dislocations, leading to secondary degenerative changes.
  \item Avascular necrosis (osteonecrosis) of the humeral head, where bone tissue dies due to an interruption of its blood supply, leading to collapse of the humeral head and severe pain.
  \item Rotator cuff tear arthropathy, a specific and severe form of arthritis associated with a massive, irreparable rotator cuff tear, primarily an indication for reverse total shoulder arthroplasty.
  \item Complex proximal humerus fractures in elderly patients with poor bone quality, where internal fixation is unlikely to achieve stable reduction or healing, or in cases of failed previous fixation.
  \item Failed previous shoulder surgeries, such as failed hemiarthroplasty (where only the humeral head was replaced) or failed open reduction and internal fixation of fractures, leading to persistent pain and dysfunction.
  \item Tumors of the proximal humerus or glenoid that necessitate resection and reconstruction of the joint.
\end{itemize}

Shoulder replacement is generally considered a primary, definitive surgical treatment for end-stage glenohumeral arthritis or other severe, irreparable shoulder conditions. In this context, it is typically pursued after a thorough and extended trial of non-operative management has failed to provide adequate pain relief or functional improvement. Non-operative strategies commonly include physical therapy, activity modification, oral anti-inflammatory medications (NSAIDs), and corticosteroid injections. When these conservative measures no longer offer symptomatic relief or functional gains, shoulder replacement serves as the primary intervention to restore joint integrity, alleviate pain, and improve overall function.

However, shoulder replacement can also serve as a secondary or salvage procedure in more complex scenarios. For instance, it may be performed as a secondary intervention following a failed hemiarthroplasty, where only the humeral head was initially replaced, and the glenoid subsequently deteriorated or developed symptomatic erosion. Similarly, it can be a crucial salvage option for complex proximal humerus fractures that are not amenable to internal fixation due to comminution or poor bone quality, or for which initial fixation has failed, leading to nonunion or malunion. The decision to proceed with shoulder replacement is always highly individualized, meticulously weighing the patient's symptoms, functional demands, overall medical health, and the precise extent and nature of the joint damage.

\section*{Relevant Surgical Anatomy}
The shoulder joint is a complex anatomical region, with the glenohumeral joint being the primary focus of shoulder replacement surgery. This synovial ball-and-socket joint is formed by the articulation of the spherical humeral head with the shallow, pear-shaped glenoid fossa of the scapula. The glenoid labrum, a fibrocartilaginous rim, deepens the glenoid cavity and enhances joint stability. Surrounding this joint are the critical rotator cuff muscles: the supraspinatus, infraspinatus, teres minor, and subscapularis. These muscles and their tendons provide dynamic stability and facilitate rotation and abduction of the arm. The deltoid muscle, a large, powerful muscle covering the shoulder, is the primary abductor of the arm and becomes particularly crucial for function after reverse total shoulder arthroplasty.

Key neurovascular structures are in close proximity to the surgical field. The axillary nerve, a branch of the posterior cord of the brachial plexus, courses inferiorly to the glenohumeral joint capsule and innervates the deltoid and teres minor muscles; it is highly susceptible to injury during glenoid exposure or humeral head dislocation. The musculocutaneous nerve and radial nerve are also in the vicinity. The axillary artery and vein, major vessels supplying and draining the upper limb, pass through the axilla and require careful protection. Lymphatic drainage of the shoulder region primarily occurs via the axillary lymph nodes. Important surgical landmarks include the coracoid process, the acromion, and the deltopectoral groove, which houses the cephalic vein and serves as the primary surgical interval for anterior approaches.

\section{Surgical Landmarks:}
\begin{itemize}
  \item Coracoid Process: A hook-like projection from the scapula, serving as an attachment point for muscles and ligaments.
  \item Acromion: The lateral extension of the scapular spine, forming the highest point of the shoulder.
  \item Deltopectoral Groove: The anatomical interval between the deltoid and pectoralis major muscles, commonly used for surgical access.
  \item Cephalic Vein: A superficial vein running within the deltopectoral groove, often retracted laterally during surgery.
\end{itemize}

\section*{Preoperative Workup \& Preparation}
Thorough preoperative workup and preparation are paramount for optimizing outcomes and minimizing potential risks associated with shoulder replacement surgery. This comprehensive process typically begins several weeks before the scheduled procedure.

Key aspects of preoperative preparation include:
\begin{itemize}
  \item **Medical Evaluation**: A comprehensive medical history and physical examination are performed by the orthopedic surgeon and often by the patient's primary care physician. This includes a full review of systems, current medications, and allergies.
  \item **Laboratory Tests**: Standard blood tests are obtained, including a complete blood count (CBC), basic metabolic panel (BMP), coagulation profile (PT/INR, PTT), and urinalysis.
  \item **Cardiac and Anesthesia Clearance**: Patients typically require clearance from their primary care physician or a cardiologist to ensure they are medically fit for major surgery and general anesthesia, often involving an electrocardiogram (ECG) and sometimes further cardiac stress testing.
  \item **Imaging Studies**: Preoperative X-rays of the shoulder are standard. A computed tomography (CT) scan is almost always obtained to precisely assess glenoid bone stock, identify any bone defects, and aid in accurate implant sizing and positioning, especially crucial for complex cases or reverse TSA. Magnetic resonance imaging (MRI) may be used to evaluate rotator cuff integrity, particularly for anatomic TSA candidates.
  \item **Medication Management**: Patients are instructed to discontinue certain medications, such as blood thinners (e.g., warfarin, clopidogrel, aspirin, NSAIDs), for a specified period before surgery to minimize the risk of intraoperative and postoperative bleeding.
  \item **NPO Status**: Patients must strictly adhere to "nil per os" (NPO) guidelines, refraining from eating or drinking for a set number of hours (typically 6-8 hours) before surgery to prevent aspiration during anesthesia induction.
  \item **Infection Screening**: Dental screening may be recommended to rule out any active dental infections, which could potentially serve as a source of bacteremia and lead to periprosthetic joint infection.
  \item **Prophylactic Antibiotics**: Intravenous antibiotics are administered shortly before the surgical incision to significantly reduce the risk of surgical site infection.
  \item **Informed Consent**: The patient receives detailed information regarding the procedure, including its potential benefits, risks, alternatives, and expected recovery course, and provides informed consent after all questions have been addressed.
\end{itemize}

**Absolute Contraindications:**
\begin{itemize}
  \item Active infection in the shoulder joint (e.g., septic arthritis) or elsewhere in the body, as this dramatically increases the risk of devastating periprosthetic joint infection.
  \item Paralysis or severe dysfunction of the deltoid muscle, which is absolutely critical for functional elevation and rotation, particularly after reverse total shoulder arthroplasty.
  \item Neuropathic joint (Charcot arthropathy), due to the rapid and progressive joint destruction and extremely high rates of implant failure.
  \item Inadequate or severely compromised bone stock in the glenoid or humerus that cannot adequately support the prosthetic components, even with bone grafting.
  \item Rapidly progressive neurological disease (e.g., advanced Parkinson's disease, amyotrophic lateral sclerosis) that would severely compromise postoperative rehabilitation or implant longevity.
  \item Uncontrolled systemic medical conditions that pose an unacceptably high anesthetic or surgical risk.
\end{itemize}

**Relative Contraindications and Precautions:**
\begin{itemize}
  \item Uncontrolled medical comorbidities, such as severe heart disease, poorly controlled diabetes mellitus, or significant pulmonary disease, which substantially increase surgical and anesthetic risks.
  \item Severe obesity (BMI >40 kg/m\textasciicircum2), which can complicate surgical exposure, increase operative time, and elevate the risk of wound complications, infection, and DVT.
  \item Active smoking, which significantly impairs wound healing, increases the risk of infection, and can compromise bone ingrowth. Patients are strongly advised to cease smoking for at least 4-6 weeks prior to surgery.
  \item Poor patient compliance or an inability to reliably participate in the rigorous postoperative rehabilitation program, which is essential for optimal functional recovery.
  \item History of prior shoulder infection, which may necessitate a two-stage revision procedure to eradicate the infection before implanting a new prosthesis.
  \item Young age with high functional demands, where concerns about implant longevity and the potential need for future revision surgery may lead to consideration of alternative treatments.
  \item Significant soft tissue deficiency or poor skin quality around the shoulder, which can compromise wound healing and increase infection risk.
\end{itemize}

\section*{Surgical Technique \& Procedure}
Shoulder replacement surgery is performed under general anesthesia, frequently supplemented with a regional nerve block, such as an interscalene block, to provide effective and prolonged postoperative pain relief. The patient is typically positioned in a semi-sitting "beach chair" position, which allows for optimal exposure and manipulation of the shoulder, or occasionally in a lateral decubitus position. The most common surgical approach is the deltopectoral incision, which involves an incision made from the coracoid process extending distally along the anterior aspect of the shoulder.

The procedure involves several key steps:
\begin{enumerate}
  \item  **Incision and Exposure**: A curvilinear deltopectoral incision is made. The deltopectoral interval, located between the deltoid and pectoralis major muscles, is identified and carefully developed. The cephalic vein is typically identified and retracted laterally to preserve it.
  \item  **Subscapularis Management**: For anatomic TSA, the subscapularis tendon, a critical component of the rotator cuff, is meticulously identified and released from its insertion on the lesser tuberosity of the humerus (subscapularis tenotomy) to gain direct access to the glenohumeral joint. For rTSA, the subscapularis may be taken down or split, depending on its quality and the surgeon's preference.
  \item  **Humeral Head Preparation**: The humeral head is dislocated from the glenoid. Precise osteotomies are performed using specialized guides to resect the damaged humeral head at the appropriate angle and depth. The humeral canal is then prepared using progressively larger reamers to create a cavity for the humeral stem component.
  \item  **Glenoid Preparation**:
      \begin{itemize}
        \item   **For Anatomic TSA**: The glenoid fossa is exposed, and any remaining articular cartilage, osteophytes, or soft tissue is removed. The glenoid surface is then reamed to a concave shape and prepared with drill holes to accept the polyethylene glenoid component, which is typically cemented into place.
        \item   **For Reverse TSA**: A glenoid baseplate is secured to the glenoid bone with multiple screws, ensuring robust fixation. A metal glenosphere (the "ball" component) is then attached to this baseplate.
      \end{itemize}
  \item  **Humeral Component Implantation**: The humeral stem is inserted into the prepared humerus, either with or without cement, depending on bone quality and implant design. For anatomic TSA, a metal humeral head component (the "ball") is then attached to the stem. For reverse TSA, a humeral tray with a polyethylene liner (the "socket") is attached to the humeral stem.
  \item  **Joint Reduction and Stability Assessment**: The prosthetic joint is carefully reduced. The surgeon then meticulously assesses joint stability throughout the range of motion, soft tissue tension, and overall biomechanics to ensure optimal function and minimize the risk of dislocation or impingement.
  \item  **Closure**: For anatomic TSA, the subscapularis tendon is meticulously repaired and reattached to its anatomical insertion site. For rTSA, the subscapularis may be repaired if feasible, or the interval closed. The surgical layers are then closed in a systematic fashion, and the skin incision is sutured or stapled. A surgical drain may be placed temporarily to manage postoperative fluid accumulation.
\end{enumerate}

\section*{Complications \& Risks}
As with any major surgical procedure, shoulder replacement carries potential risks and complications, which can be broadly categorized as general surgical risks and procedure-specific risks.

**General Surgical Risks:**
\begin{itemize}
  \item **Anesthesia Complications**: Adverse reactions to anesthetic agents, including nausea, vomiting, respiratory depression, cardiac arrhythmias, or nerve injury related to patient positioning.
  \item **Bleeding**: Excessive intraoperative blood loss potentially requiring transfusion, or postoperative hematoma formation, which may necessitate re-exploration.
  \item **Infection**: Ranging from superficial wound infection to deep periprosthetic joint infection, which can be a devastating complication often requiring multiple surgeries, including implant removal.
  \item **Deep Vein Thrombosis (DVT) and Pulmonary Embolism (PE)**: Formation of blood clots in the deep veins, typically of the lower extremities, which can dislodge and travel to the lungs, though less common in upper extremity surgery.
  \item **Cardiovascular Events**: Myocardial infarction, stroke, or other cardiac complications, particularly in patients with pre-existing cardiac disease.
\end{itemize}

**Procedure-Specific Risks:**
\begin{itemize}
  \item **Instability or Dislocation**: The prosthetic joint components may dislocate, particularly in the early postoperative period or with extreme movements, often requiring closed reduction or revision surgery.
  \item **Rotator Cuff Failure (for Anatomic TSA)**: Failure of the repaired subscapularis tendon or a new tear in other rotator cuff tendons, leading to pain and functional impairment.
  \item **Glenoid Component Loosening or Failure**: The polyethylene glenoid component can loosen from the bone, wear out over time, or fracture, necessitating revision surgery.
  \item **Humeral Component Loosening or Fracture**: The humeral stem can loosen within the humerus, or a periprosthetic fracture of the humerus can occur during or after surgery.
  \item **Nerve Injury**: Damage to nearby nerves, most commonly the axillary nerve (leading to deltoid weakness or paralysis), but also potentially the musculocutaneous nerve or components of the brachial plexus, causing widespread arm weakness or numbness.
  \item **Stiffness or Arthrofibrosis**: Formation of excessive scar tissue around the joint, severely limiting the range of motion despite diligent physical therapy.
  \item **Heterotopic Ossification**: Abnormal bone formation in the soft tissues surrounding the joint, which can restrict movement and cause pain.
  \item **Acromial or Scapular Spine Fracture (for rTSA)**: A specific risk associated with reverse total shoulder arthroplasty due to altered biomechanics and increased stress on the scapula.
  \item **Component Malposition**: Incorrect placement or alignment of the prosthetic components, leading to suboptimal function, impingement, or early implant failure.
\end{itemize}

\section*{Postoperative Care \& Recovery}
Immediately after shoulder replacement surgery, the patient is transferred to the Post-Anesthesia Care Unit (PACU) for close monitoring of vital signs, pain levels, and the neurovascular status of the operative arm. Pain management is a critical priority, often involving a multimodal approach with a combination of oral and intravenous medications, as well as the residual effects of the regional nerve block. The operative arm is typically placed in a sling or shoulder immobilizer to protect the surgical repair, limit unwanted motion, and provide comfort. Early mobilization of the hand, wrist, and elbow is encouraged to prevent stiffness and promote circulation.

Over the first 24-48 hours, the patient remains hospitalized. Physical therapy usually commences on the first postoperative day with gentle passive range of motion exercises, guided by a trained therapist, to prevent joint stiffness and initiate the healing process. Patients receive comprehensive education on proper sling use, wound care, and strict activity restrictions. Discharge typically occurs within 1-3 days, depending on the patient's pain control, mobility, and overall recovery trajectory. For the first 1-2 weeks at home, patients continue with prescribed oral pain medication, meticulous wound care, and adhere strictly to activity limitations, primarily using the sling for support and protection. A structured physical therapy program is initiated, gradually progressing from passive to active-assisted and eventually active range of motion exercises over several weeks to months, focusing on restoring strength and function while carefully protecting the healing soft tissues and implant. Full recovery can take 6-12 months.

During shoulder replacement surgery, the surgeon meticulously documents the intraoperative findings, which often include the extent of articular cartilage loss, the presence and size of osteophytes (bone spurs), the integrity and quality of the rotator cuff tendons (especially the subscapularis), and any bone defects or erosions on both the humeral head and glenoid. The condition of the surrounding soft tissues, such as the joint capsule and synovium, is also noted. These findings directly inform the choice of implant and specific surgical techniques employed. Resected tissue, primarily the humeral head and any excised osteophytes or synovial tissue, is routinely sent for pathological examination. The pathology report typically confirms the underlying diagnosis, such as severe osteoarthritis, rheumatoid arthritis, avascular necrosis, or post-traumatic arthritis, providing histological evidence of the degenerative or inflammatory processes that necessitated the surgery. This report is crucial for a complete medical record and for understanding the disease process.

A normal and successful postoperative course following shoulder replacement surgery is characterized by a predictable progression of pain resolution, functional improvement, and healing. Immediately after surgery, patients can expect moderate pain, which is effectively managed with a combination of regional blocks and oral analgesics, gradually decreasing over the first few weeks. The arm will be immobilized in a sling, and early passive range of motion exercises will begin under the guidance of a physical therapist. Over the first 6-12 weeks, patients will progressively regain passive and then active-assisted range of motion, with the goal of achieving functional movement for daily activities. Swelling and bruising around the incision site are common and typically resolve within a few weeks. The surgical incision should heal cleanly, without signs of infection. By 3-6 months, most patients experience significant pain relief and are able to perform many activities of daily living with improved comfort and mobility, although full strength and endurance may take up to a year or more of dedicated rehabilitation.

Postoperative care for shoulder replacement is a structured and prolonged process, absolutely essential for achieving optimal long-term recovery and implant longevity.
\begin{itemize}
  \item **Orthopedic Follow-up Visits**: Regular appointments with the orthopedic surgeon are scheduled, typically at 2 weeks, 6 weeks, 3 months, 6 months, and 1 year post-surgery, and then annually thereafter to monitor progress and implant status.
  \item **Suture/Staple Removal**: The surgical incision will be meticulously inspected, and sutures or staples will be removed at the first follow-up visit, usually around 10-14 days postoperatively.
  \item **Imaging**: Follow-up X-rays are routinely taken at various intervals (e.g., 6 weeks, 6 months, 1 year) to monitor implant position, assess for signs of loosening, and evaluate bone healing or adaptive remodeling.
  \item **Structured Physical Therapy**: A comprehensive and individualized physical therapy program is paramount, progressing through carefully phased stages of passive, active-assisted, and active range of motion, followed by progressive strengthening exercises. This program typically lasts for 3-6 months or longer, depending on individual progress.
  \item **Activity Resumption**: Patients receive specific guidance on gradually resuming daily activities, with strict instructions to avoid heavy lifting, repetitive overhead activities, and contact sports for an extended period (often 6-12 months or permanently for certain activities) to protect the healing tissues and implant.
\end{itemize}
Patients are educated on crucial warning signs that necessitate immediate medical attention, such as fever, chills, increasing redness or swelling around the incision, severe or worsening pain, new onset of numbness or weakness in the arm, or a sudden inability to move the arm.

\section*{Outcomes \& Prognosis}
Shoulder replacement surgery is a rapidly growing procedure globally, with incidence rates steadily increasing over the past two decades, primarily driven by an aging population, improved surgical techniques, and advancements in implant designs. It is most commonly performed in older adults, typically between 60 and 80 years of age, though it can be indicated for younger patients with specific conditions such as severe post-traumatic arthritis or avascular necrosis. There is a slight female predominance observed, particularly for primary osteoarthritis, which remains the most frequent indication for anatomic total shoulder arthroplasty. Rotator cuff tear arthropathy is the second most common indication, almost exclusively treated with reverse total shoulder arthroplasty, and its prevalence is also increasing. While the overall mortality rate associated with shoulder replacement is exceedingly low (typically less than 0.1\%), morbidity rates vary; infection rates are generally around 1-2\%, and dislocation rates range from 2-5\%, depending on the type of arthroplasty, patient comorbidities, and surgical technique.

The outcomes for shoulder replacement surgery are generally excellent, with high success rates in achieving the primary goals of significant pain relief and substantial functional improvement. For anatomic total shoulder arthroplasty (TSA), numerous studies report significant pain reduction in over 90\% of patients and good to excellent functional outcomes in 85-95\%, allowing for improved activities of daily living. Reverse total shoulder arthroplasty (rTSA) also demonstrates high rates of pain relief and often superior functional improvement for patients with rotator cuff tear arthropathy, particularly in terms of active elevation and external rotation. Implant survival rates are robust, with 10-year survival rates typically ranging from 85-95\% for both anatomic and reverse designs, and 15-year survival rates often exceeding 80\%. Recurrence of pain or loss of function can occur due to complications such as aseptic implant loosening, periprosthetic infection, or rotator cuff failure (for TSA). Prognostic factors influencing long-term success include patient age, bone quality, adherence to postoperative rehabilitation protocols, and the absence of significant medical comorbidities. Overall, shoulder replacement offers a durable and highly effective solution for end-stage shoulder pathology, significantly enhancing patients' quality of life and restoring functional independence.

\section*{References}
\begin{enumerate}
  \item MedlinePlus. Shoulder Replacement. U.S. National Library of Medicine; 2024. Available from: \url{https://medlineplus.gov/shoulderreplacement.html}
  \item Canale ST, Beaty JH. Campbell's Operative Orthopaedics. 14th ed. Philadelphia, PA: Elsevier; 2021.
  \item American Academy of Orthopaedic Surgeons (AAOS). Total Shoulder Replacement. Available from: \url{https://orthoinfo.aaos.org/en/diseases--conditions/total-shoulder-replacement/}
  \item Walch G, Boileau P, Favard L, et al. The results of reverse shoulder arthroplasty in rotator cuff deficiency. J Bone Joint Surg Am. 2005;87(12):2676-2685.
\end{enumerate}

\clearpage
